\section{Введение}
\label{sec:intro}

\subsection{Задача}
Нужно посчитать аналитическую метрику на ориентированном невзвешенном графе. Главная сложность в
том, что граф целиком не влезает в оперативную память. В условии приведён пример: памяти $128$~МБ, а
на граф нужно около $1$~ГБ. Значит, держать весь граф в ОЗУ нельзя, надо работать с диском.

Условие добавляет ещё несколько требований.
\begin{itemize}[nosep]
  \item \textbf{Обработка с диска (out-of-core).} В памяти можно держать только данные порядка числа
        вершин, но не порядка числа рёбер.
  \item \textbf{Гипер-узлы.} Встречаются вершины со степенью на порядки выше средней (средняя степень
        $50$, а у одной вершины $50\,000$). Решение не должно ломаться на таких вершинах.
  \item \textbf{Параллелизм.} Узел один, но с несколькими ядрами. Нужно занять все ядра.
  \item \textbf{Достоверность.} Результат не должен зависеть от случайности и должен повторяться от
        запуска к запуску.
  \item \textbf{Разумное время.} На больших графах счёт должен укладываться в приемлемое время.
\end{itemize}

\paragraph{Формат ввода и вывода.} Он зафиксирован условием. Вход --- CSV с заголовком, первые два
столбца \texttt{from,to}, идентификаторы вершин типа \texttt{int32}, ребро направлено
\texttt{from}~$\to$~\texttt{to}, граф невзвешенный. Выход --- CSV со столбцами \texttt{vertex,rank},
по одному числу на вершину.

\subsection{Что сделано}
Выбрана метрика \textbf{LeaderRank} --- вариант PageRank без настраиваемых параметров. Решение
написано на Java (OpenJDK~21). Рёбра графа лежат на диске и читаются потоком, в памяти живут только
массивы размера $\bigO(n)$. Счёт распараллелен по ядрам и даёт один и тот же результат при любом
числе потоков. Лимит памяти задаётся снаружи (через \texttt{-Xmx} и cgroup), а программа сама
подстраивается под него. Подробные инструкции по сборке и запуску вынесены в \texttt{README},
здесь же описано устройство решения.

\subsection{Как читать отчёт}
Раздел~\ref{sec:metric} --- про метрику LeaderRank и почему выбрана именно она.
Раздел~\ref{sec:approaches} --- варианты хранения графа и выбор одного из них.
Разделы~\ref{sec:arch}--\ref{sec:memory} --- устройство решения: архитектура, параллелизм, память.
Раздел~\ref{sec:complexity} --- анализ сложности и точности.
Раздел~\ref{sec:testing} --- тесты и замеры.
Раздел~\ref{sec:conclusion} --- выводы.
